\begin{textsample}{NEG Dim 2 – 2017\_c5 – Score -3.00 – 2017\_c5/enade\_2017\_eng\_comp\_0403c73f27.txt}  \label{ex:f2_neg_048}
CAPÍTULO 5 DISTRIBUIÇÃO DOS CONCEITOS 5.1 PANORAMA NACIONAL DA DISTRIBUIÇÃO DOS CONCEITOS O Conceito Enade é um indicador de qualidade que avalia os cursos por intermédio dos desempenhos dos estudantes, a partir dos resultados do Enade.
Os valores possíveis vão de 1 ( pior situação ) a 5 ( melhor situação ).
Os cursos com apenas um ou sem nenhum concluinte participante não obtêm o Conceito Enade, ficando Sem Conceito ( SC ).
A Tabela 5.1 apresenta a quantidade e a distribuição de cursos de Engenharia de Computação participantes do Enade/2017, por faixa de conceito e Grande Região. \textbf{Enfatiza} -se, mais uma vez, que a diferença entre os cursos tabulados neste capítulo e no capítulo 2 corresponde aos cursos Sem Conceito ( SC ), ou seja àqueles, em princípio, sem alunos concluintes que houvessem participado da prova ou àqueles que tivessem tão somente um aluno concluinte18.
Observando -se os dados da Tabela 5.1, nota -se que, dos 149 cursos participantes, 45 ( 30,2 \% ) classificaram -se com o conceito 2, conceito modal, e outros 41 cursos ( 27,5 \% ) com o conceito 3.
Nas regiões, os conceitos modais foram : 2 nas regiões Norte ( 33,3 \% ) e Sudeste ( 35,9 \% ), 2 e 4 na região Nordeste ( 30,4 \%, cada ), 4 na região Sul ( 44,0 \% ) e 3 na região Centro-Oeste ( 35,7 \% ).
O conceito 5 foi atribuído a 17 cursos do território nacional em todas as regiões : Norte ( um curso ), Nordeste ( três cursos ), Sudeste ( nove cursos ), Sul ( dois cursos ) e Centro-Oeste ( dois cursos ).
Houve, ainda, 36 cursos ( 24,2 \% ) que receberam conceito 4 e seis cursos ( 4,0 \% ) que receberam conceito 1.
Dos 149 cursos de Engenharia de Computação, quatro ( 2,7 \% ) ficaram Sem Conceito ( SC ). 18 Estes cursos com somente um concluinte não foram considerados no capítulo 2.
A região Norte participou com nove cursos ou 6,0 \% do total nacional.
Desses, três cursos receberam o conceito 2, o que equivale a 33,3 \% do total regional, além de ser o conceito modal na região, como já comentado.
A três cursos atribuíram -se os conceitos 1, 4 e 5, um para cada ( 11,1 \% ).
Foram atribuídos a outros dois ( 22,2 \% ) cursos o conceito 3.
Um curso ficou Sem Conceito ( SC ) nessa região.
A região Nordeste participou com 23 cursos ou 15,4 \% do total nacional.
Nessa região, sete cursos ( 30,4 \% em termos regionais ) obtiveram conceito 2, e outros sete cursos ficaram com conceito 4, os conceitos modais para a região.
Os conceitos 3 e 5 foram atribuídos, respectivamente, a seis e a três cursos ( 26,1 \% e 13,0 \%, respectivamente ).
Ainda nessa região, nenhum curso foi avaliado com o conceito 1 ou ficou Sem Conceito ( SC ).
Dos 78 cursos participantes da região Sudeste ( 52,3 \% do total nacional ), o conceito 2 foi o conceito modal, tendo sido obtido por 28 cursos ( 35,9 \% ).
O conceito 1 foi atribuído a cinco cursos ( 6,4 \% ), e o conceito 3, a 20 cursos ( 25,6 \% ), o segundo mais frequente.
Receberam os conceitos 4 e 5, respectivamente, 15 cursos ( 19,2 \% ) e nove cursos ( 11,5 \% ).
Nessa região, um curso ficou Sem Conceito ( SC ).
Os 25 cursos da região Sul corresponderam a 16,8 \% do total nacional.
Como já indicado, houve predominância do conceito 4, atribuído a 11 ( 44,0 \% ) dos 25 cursos da região ( conceito modal ).
O conceito 2 foi atribuído a três cursos ( 12,0 \% ), o conceito 3, a oito cursos ( 32,0 \% ), e o conceito 5, a dois cursos ( 8,0 \% ).
Um curso ficou Sem Conceito ( SC ) e nenhum curso recebeu o conceito 1 na região Sul.
Já dos 14 cursos participantes na região Centro-Oeste ( 9,4 \% do total nacional ), cinco receberam conceito 3, o conceito modal ( 35,7 \% ), como já fora destacado.
A quatro cursos ( 28,6 \% ) foi atribuído conceito 2, a dois cursos ( 14,3 \% ), conceito 4, e a outros dois cursos, conceito 5.
Nessa região, um curso ficou Sem Conceito ( SC ) e o conceito 1 não foi atribuído a nenhum curso.
O Gráfico 5.1 apresenta a distribuição cumulativa do Conceito Enade, segundo a Grande Região.
Em gráfico de distribuição cumulativa, em cada valor de abscissa ( eixo x ), o valor apresentado no eixo das ordenadas ( eixo y ) é o valor acumulado da variável.
Por exemplo, a região Norte ( linha azul ) apresenta cerca de 75,0 \% \% dos cursos com conceito 3 ou abaixo desse valor.
A região Sul ( linha roxa ), por outro lado, apresenta um pouco mais de 45 \% dos cursos na mesma situação : com conceito 3 ou abaixo.
Em linhas gerais, em um gráfico de distribuição cumulativa, poligonais mais à esquerda representam Grandes Regiões com pior distribuição, e poligonais mais à direita, Grandes Regiões com melhores conceitos.
A região Norte ( linha azul ) apresenta os cursos com conceitos mais baixos, ao passo que a região Sul ( linha roxa ) apresenta a distribuição com valores maiores, mesmo tendo a menor proporção ( quase 8 \% ) de cursos com conceito 5, quando comparada à das outras regiões.
Cursos em IES situadas nas regiões Sudeste ( linha vermelha ), Nordeste ( linha verde ) e Centro-oeste ( linha preta ) apresentam uma situação intermediária. 5.2 CONCEITOS POR CATEGORIA ADMINISTRATIVA, POR MODALIDADE DE ENSINO E GRANDE REGIÃO A Tabela 5.2 apresenta a distribuição dos cursos de Engenharia de Computação participantes do Enade/2017, por Categoria Administrativa e Modalidade de ensino, de acordo com os conceitos por eles alcançados, segundo as Grandes Regiões brasileiras.
Dos 149 cursos participantes, 58 ( 38,9 \% ) eram ministrados em Instituições Públicas e 91 ( 61,1 \% ), em Instituições Privadas.
Quanto à Modalidade de ensino, todos os cursos eram Presenciais.
De acordo com as informações da Tabela 5.2, em termos nacionais, entre Instituições Públicas e Privadas, observa -se que, dos 17 cursos avaliados com conceito 5, 14 eram oferecidos em IES Públicas e três, em IES Privadas.
Dos 58 cursos participantes de IES Públicas, o conceito 4 foi o valor modal, atribuído a 21 cursos ( 36,2 \% ).
Dentre os demais cursos de Instituições Públicas participantes, oito cursos ( 13,8 \% ) foram avaliados com conceito 2, outros 14 cursos ( 24,1 \% ) obtiveram conceito 3 e, como já comentado, 14 cursos ( 24,1 \% ), conceito 5.
Nesta categoria, um curso ficou Sem Conceito ( SC ) e nenhum curso ficou com conceito 1.
Na rede privada, o conceito modal foi 2, atribuído a 37 cursos, o correspondente a 40,7 \% dos 91 cursos da categoria.
Dentre os demais cursos participantes, seis ( 6,6 \% ) receberam conceito 1, e 27 ( 29,7 \% ), conceito 3.
O conceito 4 foi atribuído a 15 cursos ( 16,5 \% ) e o conceito 5, a três cursos ( 3,3 \% ).
Nesta Categoria Administrativa, três cursos ( 3,3 \% ) ficaram Sem Conceito ( SC ).
Como todos os cursos são da modalidade presencial, os comentários a seguir são os mesmos que já foram feitos na Tabela 5.1 para o total do Brasil.
Dos 149 cursos, 45 receberam o conceito 2 ( 30,2 \% ), conceito modal, e 41 o conceito 3 ( 27,5 \% ).
Classificaram -se com conceito 4, 36 cursos ( 24,2 \% ), 17 cursos ( 12,8 \% ) receberam conceito 5 e seis cursos ( 4,0 \% ) obtiveram conceito 1.
Quatro cursos ( 2,7 \% ) ficaram Sem Conceito ( SC ).
Na análise por região, observa -se que, na região Norte, as Instituições Privadas participaram com quatro cursos ( 44,4 \% do total regional ), dos quais nenhum ficou Sem Conceito ( SC ) ou recebeu os conceitos 4 e 5.
O conceito modal para as Instituições Privadas na região foi 2, atribuído a dois cursos, correspondendo a 50,0 \% desta categoria na região.
Dois cursos foram avaliados com conceitos 1 e 3 ( um curso em cada conceito ).
As Instituições Públicas participaram com cinco cursos na região Norte ( 55,6 \% do total regional ), e cada um deles recebeu um conceito diferente entre 2 e 5.
Um curso ficou Sem Conceito ( SC ) e nenhum curso recebeu conceito 1 nesta combinação de categoria e região.
Todos os nove cursos nesta região são presenciais, e apresentam a seguinte distribuição de conceitos : um curso recebeu conceito 1, três cursos, o conceito 2 ( modal ), dois cursos receberam o conceito 3, e dois cursos os conceitos 4 e 5 ( um cada ).
Um curso ficou Sem Conceito ( SC ).
Na região Nordeste, a rede privada participou com seis dos 23 cursos que integraram o exame, o equivalente a 26,1 \% do total da região.
Dentre os seis cursos oferecidos por IES Privadas no Nordeste, cinco cursos foram avaliados com conceito 2, o conceito modal, e um curso recebeu o conceito 4.
Nesta combinação de categoria e região, nenhum curso ficou Sem Conceito ( SC ) ou recebeu os conceitos 1, 3 e 5.
As Instituições Públicas da região Nordeste participaram com 17 cursos ( 73,9 \% do total da região, a maior proporção dessa categoria entre as regiões ).
Desses, seis ( 35,3 \% ) obtiveram conceito 3 e outros seis, conceito 4, os valores modais.
Seguidos por outros três cursos ( 17,6 \% ), que receberam conceito 5 e, dois cursos que receberam conceito 2.
Nenhum curso ficou Sem Conceito ( SC ) ou com conceito 1.
Também nessa região, todos os cursos foram presenciais : 23.
As modas foram os conceitos 2 e 4, com sete cursos cada ( 30,4 \% ).
Os demais cursos presenciais da região foram alocados aos conceitos 3 ( seis cursos, 26,1 \% ) e 5 ( três cursos, 13,0 \% ).
Nenhum curso ficou sem ( SC ) ou com conceito 1.
Na região Sudeste, a proporção de cursos da rede privada, 73,1 \%, foi mais \textbf{elevada} do que a registrada nas demais regiões brasileiras, correspondendo a 57 dos 78 cursos participantes.
Nesta categoria e região, o conceito modal foi 2, atribuído a 25 cursos ( 43,9 \% ).
Um curso ficou Sem Conceito ( SC ).
Os demais foram avaliados com conceito 1 ( cinco cursos, 8,8 \% ), conceito 3 ( 16 cursos, 28,1 \% ), conceito 4 ( oito cursos, 14,0 \% ) e conceito 5 ( dois cursos, 3,5 \% ).
Dentre os 21 cursos oferecidos por Instituições Públicas na região Sudeste, os conceitos modais foram 4 e 5, atribuído a sete cursos cada ( 33,3 \% ).
Os demais cursos receberam os conceitos 2 e 3, respectivamente, três e quatro cursos.
Na região Sudeste, nesta categoria, nenhum curso ficou Sem Conceito ( SC ) ou foi alocado ao conceito 1.
Também na região Sudeste, todos os cursos foram presenciais.
A moda foi o conceito 2, com 28 cursos ( 35,9 \% ).
Os demais cursos presenciais da região foram alocados aos conceitos 1 ( cinco cursos, 6,4 \% ), 3 ( 20 cursos, 25,6 \% ), 4 ( 15 cursos, 19,5 \% ) e 5 ( nove cursos, 11,5 \% ).
Apenas um curso ficou Sem Conceito ( SC ).
As Instituições Privadas concentraram 14 dos 25 cursos participantes da região Sul, 56,0 \% do total regional.
Desses, seis ( 42,9 \% ) obtiveram conceito 4, o conceito modal.
Nesta combinação de Categoria Administrativa e Grande Região, os conceitos 2 e 3 foram atribuídos a dois ( 14,3 \% ) e a cinco ( 35,7 \% ) cursos, respectivamente.
Nenhum curso recebeu os conceitos 5 ou 1 e apenas um ficou Sem Conceito ( SC ).
As Instituições Públicas na região Sul participaram com 11 cursos, 44,0 \% dos cursos da região.
O conceito modal foi 4 ( 45,5 \% ), atribuído a cinco cursos.
Aos demais cursos foram atribuídos os conceitos 2 ( um curso, 9,1 \% ), 3 ( três cursos, 27,3 \% ) e 5 ( dois cursos, 18,2 \% ).
Nenhum curso da região Sul oferecido por IES Pública ficou Sem Conceito ( SC ) ou recebeu conceito 1.
Também na região Sul, todos os cursos foram presenciais.
A moda foi o conceito 4, com 11 cursos ( 44,0 \% ).
Os demais cursos presenciais da região foram alocados aos conceitos 2 ( três cursos, 12,0 \% ), 3 ( oito cursos, 32,0 \% ) e 5 ( dois cursos, 8,0 \% ).
Somente um curso ficou Sem Conceito ( SC ) e o conceito 1 não foi atribuído a nenhum curso.
Na região Centro-Oeste, dez dos 14 cursos participantes eram de Instituições Privadas ( 71,4 \% em termos regionais ).
Desses, cinco ( 50,0 \% ) receberam conceito 3, o conceito modal.
Dentre os demais, um recebeu conceito 5 ( 10,0 \% ) e três receberam conceito 2 ( 30,0 \% ).
Nenhum curso foi alocado no conceito 1 ou 4, e um curso ficou Sem Conceito ( SC ).
Quanto aos quatro cursos oferecidos por Instituições Públicas na região Centro-Oeste ( 28,6 \% do total regional ), dois cursos ( 50,0 \% ) apresentaram conceito 4, conceito modal, e dois ficaram com conceitos 2 e 5 ( um curso, cada ).
Nenhum curso ficou Sem Conceito ( SC ) ou com conceitos 1 ou 3.
Nessa região, todos os cursos foram presenciais, e o conceito modal foi o 3, sendo atribuído a cinco cursos ( 35,7 \% ).
Quatro cursos obtiveram o conceito 2 ( 28,6 \% ), dois cursos, o conceito 4 ( 21,4 \% ) e outros dois cursos, conceito 5.
Dois cursos receberam o conceito 5 ( 14,3 \% ).
Um curso ficou Sem Conceito ( SC ) e nenhum recebeu conceito 1.
O Gráfico 5.2 apresenta a distribuição do Conceito Enade segundo a Categoria Administrativa da IES.
Os cursos em IES Públicas ( linha azul ) apresentam uma poligonal mais à direita do que os em IES Privadas ( linha verde ), onde as porcentagens acumuladas dos cursos com conceito 3 ou abaixo são de, respectivamente, 37,9 \% e 78,4 \%, o que indica que cursos em IES Públicas apresentam valores de conceito maiores. 5.3 CONCEITOS POR ORGANIZAÇÃO ACADÊMICA E GRANDE REGIÃO Na Tabela 5.3, encontra -se a distribuição dos conceitos atribuídos aos cursos participantes do Enade/2017 na Área de Engenharia de Computação, por Organização Acadêmica, segundo as Grandes Regiões brasileiras.
Dos 149 cursos de Engenharia de Computação participantes, 91 eram oferecidos em Universidades ; 29, em Centros Universitários ; 24, em Faculdades ; e 5, em CEFET/IFET.
Essa distribuição corresponde a, respectivamente, 61,1 \%, 19,5 \%, 16,1 \% e 3,4 \% dos cursos.
De acordo com os dados apresentados, de todos os 17 cursos avaliados com conceito 5, 13 eram vinculados a Universidades.
Os cursos vinculados a esse tipo de Organização Acadêmica tiveram o conceito 4 como conceito modal, atribuído a 26 cursos ( 28,6 \% ).
Os demais cursos vinculados a Universidades receberam os conceitos 1 ( três cursos, 3,3 \% ), 2 ( 22 cursos, 24,2 \% ), 3 ( 23 cursos, 25,3 \% ) e 5 ( 13 cursos, como já mencionado, 14,3 \% ).
Ficaram Sem Conceito ( SC ), quatro cursos ( 4,4 \% ).
Dentre os cursos em Centros Universitários, o conceito modal foi 3, atribuído a 11 cursos ( 37,9 \% ).
Nenhum curso vinculado a esse tipo de Organização Acadêmica ficou Sem Conceito ( SC ).
Os outros cursos receberam os conceitos 1 ( três cursos, 10,3 \% ), 2 ( nove cursos, 31,0 \% ), 4 ( cinco cursos, 17,2 \% ) e 5 ( um curso, 3,4 \% ).
Dos 24 cursos mantidos por Faculdades, 13 ( 54,2 \% ) receberam conceito 2, o conceito modal.
O segundo conceito mais frequentemente atribuído aos cursos oferecidos por essa categoria de IES foi o conceito 3, recebido por seis cursos ( 25,0 \% ).
Na sequência, três cursos obtiveram conceito 4, ( 12,5 \% ) e dois cursos ( 8,3 \% ), conceito 5.
Nenhum curso foi atribuído ao conceito 1 ou ficou Sem Conceito ( SC ).
Dos cinco cursos mantidos por CEFET/IFET, dois ( 40,0 \% ) receberam conceito 4, o conceito modal.
Os outros três cursos receberam os conceitos 2, 3 e 5 ( um curso, cada ).
Nenhum curso ficou Sem Conceito ( SC ) ou recebeu o conceito 1.
Considerando -se separadamente as regiões brasileiras, verifica -se que, na região Norte, as Universidades integraram seis dos nove cursos participantes.
Todos eles \textbf{distribuídos} entre os conceitos 1 e 5, com um curso cada.
Além disso, um curso ficou Sem Conceito ( SC ).
Os Centros Universitários da região Norte foram representados por um curso, que recebeu conceito 3.
As Faculdades participaram com dois cursos na região Norte, os quais receberam conceito 2.
Os CEFET/IFET não registraram cursos nessa região.
Na região Nordeste, as Universidades participaram com 19 ( 82,6 \% ) dos 23 cursos da Área de Engenharia de Computação oferecidos na região, dos quais nenhum ficou Sem Conceito ( SC ) ou recebeu conceito 1.
Nota -se que tais cursos apresentaram uma distribuição com modas \textbf{distribuídas} aos conceitos 3 e 4, com seis cursos cada ( 31,6 \% ).
Obtiveram os conceitos 2 e 5, respectivamente, quatro ( 21,1 \% ) e três cursos ( 15,8 \% ).
Os Centros Universitários contaram com um curso participante na região Nordeste, que obteve conceito 2.
As Faculdades foram representadas por dois cursos na região Nordeste, que também receberam conceito 2.
Os CEFET/IFET participaram com um curso, o qual recebeu conceito 4.
Na região Sudeste, as Universidades concentraram 37 ( 47,4 \% ) dos 78 cursos de Engenharia de Computação da região.
Dentre os cursos em Universidades, nessa região, o conceito modal foi 2, atribuído a 13 cursos ( 35,1 \% ).
Os demais cursos receberam os conceitos 1 ( dois cursos, 5,4 \% ), 3 ( oito cursos, 21,6 \% ), 4 ( sete cursos, 18,9 \% ) e 5 ( seis cursos, 16,2 \% ).
Um curso ficou Sem Conceito ( SC ).
Os Centros Universitários participaram com 20 cursos na região Sudeste ( 25,6 \% ), dos quais seis ( 30,0 \% ) obtiveram o conceito 2 e outros seis, conceito 3, os conceitos modais.
O conceito 1 foi atribuído a três cursos ( 15,0 \% ), o conceito 4, a quatro cursos ( 20,0 \% ) e, por fim, o conceito 5, a um curso ( 5,0 \% ).
Nenhum curso ficou Sem Conceito ( SC ) nesta combinação de organização acadêmica e região.
Já as Faculdades foram representadas por 18 cursos na região Sudeste ( 23,1 \% ), dos quais nenhum ficou Sem Conceito ( SC ) ou recebeu conceito 1.
Aos demais cursos foram atribuídos os conceitos 2 ( nove cursos, o conceito modal, 50,0 \% ), 3 ( cinco cursos, 27,8 \% ), 4 ( três cursos, 16,7 \% ) e 5 ( um curso, 5,6 \% ).
Os CEFET/IFET participaram com três, os quais receberam conceitos 3, 4 e 5, um curso cada.
Dos 25 cursos da região Sul, 22 ( 88,0 \% ) eram de Universidades, o maior percentual dessa categoria entre as regiões, para os quais o conceito modal foi 4, atribuído a dez cursos ( 45,5 \% ).
Nesse tipo de organização, um curso ficou Sem Conceito ( SC ) e nenhum curso recebeu conceito 1.
Os demais cursos receberam os conceitos 2 ( três cursos, 13,6 \% ), 3 ( seis cursos, 27,3 \% ) e 5 ( dois cursos, 9,1 \% ).
Dos dois cursos participantes de Centros Universitários da região Sul, um recebeu o conceito 3 e o outro o conceito 4.
O único curso vinculado a Faculdades na região Sul, recebeu conceito 3.
Nesta região, não foram avaliados cursos em CEFET/IFET.
Na região Centro-Oeste, sete dos 14 cursos eram de Universidades ( 50,0 \% ).
Nesse tipo de organização, um curso ficou Sem Conceito ( SC ).
Os conceitos 2 e 5 foram atribuídos a um curso cada, e os conceito 3 e 4, a dois cursos cada, os conceitos modais.
O conceito 1 não foi atribuído a nenhum curso.
Os Centros Universitários da região Centro-Oeste contaram com cinco cursos, que obtiveram os conceitos 3 ( três cursos, conceito modal ) e 2 ( dois cursos ).
O curso oferecido em Faculdades na região Centro-Oeste recebeu conceito 5, o conceito máximo.
Nessa região, um curso foi avaliado em CEFET/IFET e recebeu conceito 2.
O Gráfico 5.3 apresenta a distribuição do Conceito Enade segundo a Organização Acadêmica da IES.
Os cursos em CEFET/IFET ( linha roxa ) apresentam uma poligonal mais à direita do que as demais, denotando uma melhor distribuição de Conceitos Enade.
Nota -se também que nenhum dos cursos em CEFET/IFET ( linha roxa ) obtiveram conceitos 1, já que a poligonal cumulativa parte de 0 \% no conceito 2.
No outro extremo, estão os cursos oferecidos em Faculdades ( linha vermelha ) e Centros Universitários ( linha verde ) com as poligonais mais à esquerda e, consequentemente, uma distribuição de Conceitos Enade com menores valores, diferenciando -se em dois aspectos : apenas os cursos em IES de Centros Universitários ( linha verde ) obtiveram conceito 1 ( 6,9 \% ), e no conceito 2 ou abaixo, houve maior porcentagem acumulada de cursos em IES de Faculdades ( 54,2 \% contra 34,5 \% dos de IES de Centros Universitários ) ( linha vermelha ).
IES de Universidades ( linha azul ) apresentam uma situação intermediária.

% matched lemmas: distribuir, elevado, enfatizar
\end{textsample}
